\section{Discussion \& Ablation Analysis}
\label{sec:discussion}

\paragraph{Method Lineage as Progressive Macro-Ablations.}
Our development progression provides a systematic, step-wise ablation of the architectural components culminating in THIEF:
\begin{itemize}[leftmargin=*,nosep]
    \item \textbf{Monolithic vs.\ Modular Capacity (MARC vs.\ COOP):} Transitioning from a single monolithic policy with potential-based credit shaping (MARC) to a 2-expert value-bidding architecture (COOP) elevates Stage 3 win rate from $31.4\%$ to $35.5\%$, and Stage 4 escape rate from $6.5\%$ to $7.8\%$. However, COOP exhibits severe routing imbalance: without load balancing, a single expert captures over $92\%$ of rollout dispatches, reducing the multi-expert capacity back toward monolithic behavior.
    \item \textbf{Static vs.\ Evolutionary Recombination (COOP vs.\ E-COOP):} Incorporating Fisher Information parameter recombination in E-COOP enables rapid policy adaptation across stages, matching COOP on Stage 4 ($7.6\% \pm 1.4\%$). However, E-COOP consolidates all adapted parameters into a single champion network ($K=1$). By overwriting historical parameters during consolidation, E-COOP loses the ability to retain distinct specialized policies for disparate subtasks.
    \item \textbf{Consolidation vs.\ Dynamic Multi-Expert Pools (E-COOP vs.\ THIEF):} E-COOP functions as the primary macro-architectural ablation of THIEF: it employs Fisher recombination but ablates dynamic spawning, sandbox warmup incubation, and load-balancing regularization. By dynamically expanding the expert pool on detected plateaus and preserving specialists via hysteresis barriers and load-balancing losses, THIEF elevates Stage 4 win rate from $7.6\%$ to \textbf{$13.3\% \pm 2.4\%$} (a $+75.0\%$ relative gain), proving that retaining specialized modular sub-networks is essential for long-horizon spatial coordination.
\end{itemize}

\paragraph{Component Micro-Ablation Breakdown.}
To evaluate each internal mechanism of THIEF in controlled isolation, we benchmark standalone ablations on diagnostic stages (Stage 2 and Stage 4; Table~\ref{tab:ablations} in Appendix):
\begin{itemize}[leftmargin=*,nosep]
    \item \textbf{Impact of Sandbox Warmup Incubation:} Disabling isolated sandbox warmup ($\tau_{\text{iso}}=0$) reduces Stage 4 win rate from $13.7\%$ to $12.4\%$. Without isolated incubation, newly spawned specialists suffer from uncalibrated critics, causing them to be immediately outbid by established generalists in global dispatch.
    \item \textbf{Impact of Curvature-Optimal Recombination:} Replacing Fisher Information recombination with single-parent cloning with noise (\texttt{clone\_best}) degrades Stage 4 win rate to $12.4\%$ and inflates the expert switch rate to $7.5\%$. Combining parent parameters along curvature manifolds preserves established coordination primitives while exploring orthogonal subtask directions.
    \item \textbf{Impact of Load-Balancing Regularization:} Removing the load-balancing entropy loss ($\alpha_{\text{bal}}=0$) lowers Stage 4 win rate to $12.6\% \pm 2.3\%$ and produces textbook routing monopoly: in every run that spawned additional experts, the empirical routing entropy collapsed to $H(f) = 0.00$ and the Gini expert imbalance rose from $0.33$ (full THIEF) to $0.50$, confirming Proposition~\ref{prop:entropy}'s anti-monopoly safeguard empirically.
    \item \textbf{Impact of Hysteresis Routing:} Eliminating the switching barrier ($\epsilon=0$) removes the mechanism's own anti-thrashing telemetry, so the measured switch rate no longer isolates policy churn: the remaining switching occurs almost exclusively during spawn-driven exploration windows (Table~\ref{tab:ablations} shows $2.8\%$ aggregate, versus $4.7\%$ for full THIEF, whose hysteresis makes routine step-to-step switching free). The clearest cost appears in downstream stability: the no-hysteresis variant shows the widest spread between Stage 2 and Stage 4 held-out performance among non-degraded variants, consistent with the policy-thrashing pathology the barrier is designed to prevent.
    \item \textbf{Impact of Plateau Triggering:} Replacing empirical win-rate plateau triggering with a fixed update schedule produces the largest Stage 4 degradation ($10.3\% \pm 0.9\%$) and indiscriminate spawning ($3.0$ spawns and a $25.2\%$ switch rate per run versus $0.5$ and $4.7\%$ for full THIEF), demonstrating that deficit-driven expansion is far more sample-efficient than arbitrary periodic spawning.
\end{itemize}

\paragraph{Computational Complexity \& Wall-Clock Overhead.}
THIEF introduces minimal computational overhead relative to standard PPO. The empirical Fisher Information Matrix is maintained as a running diagonal vector requiring $\mathcal{O}(D)$ memory per expert ($D \approx 1.8 \times 10^5$ parameters). Parameter recombination and deficit gradient steps occur exclusively during plateau events, incurring an $\mathcal{O}(KD)$ parameter operation ($K \le 8$ active experts) that completes in less than $0.15$ seconds on GPU. Over the entire $3.4 \times 10^6$ step training curriculum, spawning operations account for less than $0.3\%$ of total wall-clock execution time.

\paragraph{Limitations \& Future Directions.}
While THIEF consistently achieves the highest win rate on every stage, absolute performance on the largest benchmark remains modest for all methods: on the $50 \times 50$ Stage 4 facility, THIEF succeeds in $13.3\%$ of episodes and the strongest baselines in $1.4$--$9.4\%$, i.e., every evaluated method fails the overwhelming majority of the time. Our failure analysis (Section~\ref{sec:results}) shows that $70.2\%$ of THIEF's unsuccessful episodes are caused by alarm accumulation during the extraction phase under $7 \times 7$ egocentric partial observability: agents cannot observe global guard movements outside their local window, and deep long-horizon coordination---rather than subtask representation, where THIEF already dominates---is the binding constraint at this scale. We therefore position Stage 4 as a stress test that the entire field currently fails, with THIEF degrading slowest and holding the largest margin over the second-best method (ROMA at $9.4\%$). Future work will investigate integrating graph neural networks for peer-to-peer communicative coordination, hierarchical macro-action abstractions, and partial-observability-aware alarm prediction to address this bottleneck.
